\teachcover{Continued fractions\\and matrix fixed points}{1}{
One repeating block, one quadratic equation
\[
  x=\frac{ax+b}{cx+d}
\]
}{
Open with the question: why should repeating continued fractions know
anything about quadratic equations? Students have already met linear
fractional transformations. Treat the first few slides as a review,
not a new chapter of prerequisites. Today's destination is the
prefix-conjugation argument. Source: course book, Sections 3.2 and 3.7.
}

\teachslide{A repeating pattern}{2}{
\[
 \sqrt2=[1;2,2,2,\ldots]=[1;\overline2]
\]
\[
 1+\sqrt2=[2;2,2,2,\ldots]=[\overline2]
\]
\spacer
The first expansion is \key{eventually periodic}.

The second is \key{purely periodic}.
}{
Ask students which digit belongs to the repeating block in each
expression. Explain that pure periodicity includes the integer part.
A bar over the block means to repeat it indefinitely. The original
number and its complete quotient can have different kinds of
periodicity. Recall that rational simple continued fractions terminate.
In these lectures, a finite rational expansion does not count as
eventually periodic. Source: course book, Section 3.7.
}

\teachslide{The two periodicity theorems}{2}{
A \key{quadratic irrational} is an irrational real root of an integer quadratic.
\spacer
\textbf{Lagrange}
\[
 \text{eventually periodic}
 \quad\Longleftrightarrow\quad
 \text{quadratic irrational}
\]
\spacer
\textbf{Galois}

A quadratic irrational $x$ is purely periodic exactly when
\[
 \key{x>1\qquad\text{and}\qquad -1<x'<0.}
\]
Here $x'$ is the other root of its quadratic equation.
}{
State the destination, without trying to prove both results now.
Use $\sqrt2$ as the simplest example of a quadratic irrational.
Today proves the forward direction of Lagrange. Lecture 2 proves the
converse and both directions of Galois. Clarify the names gently:
Legendre's continued-fraction theorem concerns exceptionally good
rational approximations. Lagrange is the eventual-periodicity theorem.
The theorem statements use simple continued fractions. Source:
course book, Sections 3.7--3.9. External cross-check: Corless,
\href{https://www.cecm.sfu.ca/organics/papers/corless/confrac/html/node6-an.shtml}{Periodic and Fixed Points of the Gauss Map}.
}

\teachslide{A M\"obius transformation}{2}{
\[
 M=\mat abcd,\qquad
 \key{T_M(t)=\frac{at+b}{ct+d}}
\]
\[
 a,b,c,d\in\Q,\qquad ad-bc\ne0
\]
\spacer
We use the formula where $ct+d\ne0$.

The determinant condition makes the transformation invertible.
}{
Use the notation from your preceding lecture, while making clear that
$T_M(t)$ is a fractional action rather than ordinary matrix
multiplication by a scalar. The denominator condition is essential.
With rational coefficients, any finite pole is rational, so an
irrational input never hits a pole. Infinity is a convenient optional
convention, not a prerequisite for today's proofs. Source: book 3.2.
}

\teachslide{Why a matrix acts on a number}{2}{
\[
 \mat abcd\begin{pmatrix}t\\1\end{pmatrix}
   =\begin{pmatrix}at+b\\ct+d\end{pmatrix}
\]
\spacer
Take the ratio of the two output coordinates.
\[
 t\ \longmapsto\ \frac{at+b}{ct+d}
\]
\spacer
Multiplying every entry of $M$ by the same nonzero number
leaves $T_M$ unchanged.
}{
This is the projective viewpoint without projective-space terminology.
Emphasize the ratio. The columns are ordinary vectors, and the
fractional transformation records their coordinate ratio.
Ask why $M$ and $-M$ produce the same map. The identity matrix and
all its nonzero scalar multiples therefore represent the identity
transformation. We will need that observation when excluding
degenerate fixed-point equations. Source: book 3.2 and Appendix A.2.
}

\teachslide{Composition}{3}{
\[
 \key{T_M\circ T_N=T_{MN}}
\]
\spacer
The rightmost matrix acts first.
\[
 M=\mat1101,\qquad N=\mat0110
\]
\[
 T_M(T_N(t))=1+\frac1t,\qquad
 MN=\mat1110
\]
}{
Here $M$ is translation by 1 and $N$ is inversion. Compute the
product with the class, then compare it with the nested function.
If students already remember the general substitution proof,
do not repeat its four coefficients. The first backup slide
contains that calculation if requested. Ask what happens in the
opposite order: $NM=\mat0111$ acts by $1/(t+1)$.
This makes the order of multiplication memorable. Source: book 3.2.
}

\teachslide{An inverse, by solving for $t$}{2}{
\[
 y=\frac{at+b}{ct+d}
 \qquad\Longrightarrow\qquad
 \key{t=\frac{dy-b}{a-cy}}
\]
\spacer
\[
 M^{-1}=\frac1{ad-bc}\mat d{-b}{-c}a
\]
\[
 (T_M)^{-1}=T_{M^{-1}}
\]
The scalar factor disappears when we take a coordinate ratio.
}{
Let students supply the rearrangement
$(a-cy)t=dy-b$. An important corollary is that an invertible
rational M\"obius transformation preserves irrationality:
if its output were rational, its inverse would make the input
rational. There is no hidden pole in this argument.
For $y=T_M(t)$, $a-cy=(ad-bc)/(ct+d)\ne0$.
We use this fact later to rule out a merely linear equation.
Source: book 3.2 and the proof in 3.7.
}

\teachslide{One continued-fraction digit}{2}{
\[
 \key{D(a)=\mat a110}
\]
\[
 T_{D(a)}(t)=a+\frac1t
\]
\spacer
The matrix inserts one digit before the remaining tail.
\[
 \det D(a)=-1,\qquad D(a)^{-1}=\mat011{-a}
\]
}{
The tail $t$ here is a whole remaining continued fraction.
The letter $D$ in $D(a)$ means a digit matrix. In Lecture 2,
the plain letter $D$ will instead denote a discriminant under
a square root; point out the contextual distinction if needed.
All finite digit products have determinant $\pm1$ and integer
inverses. This is the exact sense in which we can undo a prefix.
Source: book 3.2--3.3.
}

\teachslide{A block of digits}{2}{
\[
 [a_0;a_1,\ldots,a_{r-1},t]
 =T_{D(a_0)D(a_1)\cdots D(a_{r-1})}(t)
\]
\spacer
For the block $1,2$,
\[
 D(1)D(2)=\mat3121
\]
\[
 1+\frac1{2+1/t}=\frac{3t+1}{2t+1}.
\]
}{
Translate the displayed continued fraction back into nested
functions. All digits after the integer part are positive.
The last entry $t$ is the tail, so it need not be an integer.
The example fixes the period matrix that will run through the
rest of the lecture. Have students keep $\mat3121$ on their
paper. Source: book 3.2 and the worked example in 3.7.
}

\teachslide{A repeating tail}{2}{
\[
 x=[\overline{b_0,\ldots,b_{\ell-1}}],
 \qquad b_i\ge1
\]
\[
 N=D(b_0)\cdots D(b_{\ell-1})
   =\mat ABC{D_0}
\]
\spacer
Removing one whole period leaves the same tail.
\[
 \key{x=T_N(x)=\frac{Ax+B}{Cx+D_0}}
\]
Here $B,C>0$ and $\det N=(-1)^\ell$.
}{
This is the central idea. Pause before showing the fixed-point
equation and ask what remains after removing a complete block.
An infinite simple continued fraction has a well-defined
irrational limit; this follows from the convergent bounds
already developed in class. A backup gives the short
irrationality argument. Positivity of the period digits makes
$B,C$ strictly positive, even for a one-digit period.
Source: book 3.7.
}

\teachslide{A fixed point gives an equation}{2}{
\[
 x=\frac{Ax+B}{Cx+D_0}
\]
\[
 Cx^2+D_0x=Ax+B
\]
\[
 \key{Cx^2+(D_0-A)x-B=0}
\]
\spacer
Since $C>0$, this is a genuine quadratic equation.

Thus a purely periodic simple continued fraction is quadratic.
}{
Keep this algebraic step brief.
The nontrivial part will be understanding exactly which
quadratics repeat. At this point the number is irrational
because the continued fraction is infinite, so it is a
quadratic irrational, not merely some root of a reducible
quadratic. Contrast the period matrix with the identity:
$t=T_I(t)$ only yields $0=0$. Source: book 3.7.
}

\teachslide{Your turn: the block $1,2$}{3}{
\[
 x=[\overline{1,2}],\qquad N=\mat3121
\]
\spacer
Find a quadratic equation for $x$.

Find both roots.

Which root can be the continued fraction?
\spacer
\key{Take one minute before we compare answers.}
}{
Give students a real minute of silence, then ask for a
quadratic before asking for its roots. Expected equation:
$2x^2-2x-1=0$. Expected roots: $(1\pm\sqrt3)/2$.
The continued fraction exceeds 1, so it is the positive root.
If someone takes the negative root, ask which interval contains
$1+1/(2+1/(\cdots))$. Advance only after discussion.
This is book Exercise 15, expanded into a classroom prompt.
}

\teachslide{The periodic root}{2}{
\[
 2x^2-2x-1=0
\]
\[
 \key{x=[\overline{1,2}]=\frac{1+\sqrt3}{2}}
\]
\spacer
The other root is
\[
 x'=\frac{1-\sqrt3}{2},\qquad -1<x'<0.
\]
}{
The prime denotes the other algebraic root, not a derivative.
Do not yet explain the significance of the interval $(-1,0)$.
Point out the pattern and promise that Lecture 2 will
explain it exactly. One can check the value by substituting
it into $(3x+1)/(2x+1)$. Source: book 3.7.
}

\teachslide{A finite prefix}{2}{
\[
 x=T_N(x),\qquad y=T_M(x)
\]
\spacer
The original number $y$ may have extra digits before the period.
\[
 y=[1;\overline{1,2}]=1+\frac1x,
 \qquad M=D(1).
\]
\spacer
Which transformation fixes \key{$y$}?
}{
Emphasize the distinction between a periodic tail and the original
number. We know a quadratic equation for the tail
$x$, but must produce one for the number $y$ with its prefix.
Let students suggest substitution or inversion.
Direct substitution works, but conjugation packages it into
one reusable identity. Source: book 3.7, equation (3.3).
}

\teachslide{The prefix conjugation identity}{3}{
\[
 \begin{aligned}
 y&=T_M(x)\\[3pt]
  &=T_M\bigl(T_N(x)\bigr)\\[3pt]
  &=T_M\bigl(T_N(T_{M^{-1}}(y))\bigr).
 \end{aligned}
\]
\spacer
\[
 \key{y=T_{MNM^{-1}}(y)}
\]
Attaching a prefix conjugates the transformation that fixes the tail.
}{
Read the composition from right to left: remove the prefix,
apply the period to its fixed tail, and put the prefix back.
Spend the full three minutes on why each equality holds.
For students comfortable with function notation, write
$MNM^{-1}$ directly while saying that the matrices act
fractionally. An empty prefix means $M=I$.
Source: book 3.7, equation (3.3).
}

\teachslide{Why the new equation is genuinely quadratic}{3}{
Put $K=MNM^{-1}$.

\spacer
$K$ has integer entries and determinant $\pm1$.

$K$ is nonscalar, since conjugating a scalar matrix gives a scalar matrix.

$y$ is irrational, since $T_{M^{-1}}(y)=x$.
\spacer
For $K=\mat abcd$,
\[
 \key{cy^2+(d-a)y-b=0.}
\]
This polynomial is nonzero and cannot be merely linear.
}{
This is the rigor checkpoint. If $K=\lambda I$, then
$N=M^{-1}KM=\lambda I$, contradicting the positive
lower-left entry of the period matrix. Thus the three
coefficients $c,d-a,-b$ are not all zero.
If $c=0$, a nonzero rational linear polynomial would vanish
at the irrational $y$, which is impossible.
Rationality of a M\"obius image would also imply rationality
of its inverse image. The backup fixed-point lemma gives
this argument in isolation. Source: book 3.7.
}

\teachslide{The equation for $\sqrt3$}{3}{
\[
 M=\mat1110,\qquad
 N=\mat3121,\qquad
 M^{-1}=\mat011{-1}
\]
\[
 MNM^{-1}=\mat2312
\]
\[
 y=\frac{2y+3}{y+2}
 \qquad\Longrightarrow\qquad y^2=3
\]
\[
 \key{[1;\overline{1,2}]=\sqrt3}
\]
}{
Have students check at least one row of the matrix product.
For an intermediate check, $MN=\mat5231$.
The positive sign follows from $y=1+1/x>0$.
This example both verifies the familiar continued fraction
of $\sqrt3$ and shows what conjugation actually does to
the quadratic equation. Source: book 3.7 and Exercise 21.
}

\teachslide{Your turn: a different prefix}{3}{
Keep $x=[\overline{1,2}]$, but use $M=D(2)$.
\[
 y=[2;\overline{1,2}]=2+\frac1x
\]
\spacer
Compute $MNM^{-1}$ and identify $y$.

Can you write its digits with a bar starting at the first digit?
\spacer
\key{A written prefix need not be a genuine preperiod.}
}{
Allow about 90 seconds for the calculation.
Students who know $1/x=\sqrt3-1$ can immediately identify
$y=1+\sqrt3$, but ask them to check the conjugated matrix too.
The result is $\mat3211$, giving $y^2-2y-2=0$.
The digit stream is $2,1,2,1,\ldots$, so its pure period
is $[\overline{2,1}]$. Source: book Exercise 21.
}

\teachslide{A prefix that disappears}{2}{
\[
 D(2)ND(2)^{-1}=\mat3211
\]
\[
 y^2-2y-2=0,\qquad y=1+\sqrt3
\]
\spacer
\[
 \key{[2;\overline{1,2}]=[\overline{2,1}]}
\]
\spacer
The digits repeat from the beginning after all.
}{
Distinguish a particular written decomposition into prefix
and period from the shortest possible preperiod of the
actual digit sequence. Here the first digit already belongs
to a cycle of length two. This is an important example to
bring back after Galois's criterion, since the conjugate
$1-\sqrt3$ lies in $(-1,0)$. Source: book 3.9.
}

\teachslide{Quadratic conjugates}{2}{
\[
 x=u+v\sqrt D,\qquad
 \key{x'=u-v\sqrt D}\qquad (u,v\in\Q,\ v\ne0)
\]
\spacer
Replacing $\sqrt D$ by $-\sqrt D$ respects rational arithmetic.
\[
 \bigl(T_M(x)\bigr)'=T_M(x')
\]
\spacer
For example,
\[
 (\sqrt2)'=-\sqrt2,\qquad
 (1+\sqrt2)'=1-\sqrt2.
\]
}{
Assume $D$ is a positive nonsquare integer.
Clarify again that the prime is algebraic conjugation,
not differentiation. All we need is that changing the
sign of $\sqrt D$ commutes with adding, multiplying,
and taking reciprocals when the denominators are nonzero.
The complete quotients are rational transformations of
the original surd, so their conjugates follow the same
algebra. Source: book 3.7--3.8.
}

\teachslide{Reduced surds}{3}{
A quadratic irrational is \key{reduced} when
\[
 \key{x>1\qquad\text{and}\qquad -1<x'<0.}
\]
\spacer
\begin{center}
\renewcommand{\arraystretch}{1.25}
\begin{tabular}{lll}
\toprule
$x$ & $x'$ & Reduced?\\
\midrule
$\sqrt2$ & $-\sqrt2$ & No\\
$1+\sqrt2$ & $1-\sqrt2$ & Yes\\
$\sqrt2/2$ & $-\sqrt2/2$ & No\\
\bottomrule
\end{tabular}
\end{center}
}{
Ask which inequality fails in each of the two negative
examples. For $\sqrt2$ it is the lower bound on the
conjugate. For $\sqrt2/2$ it is $x>1$.
Galois says that this definition is exactly the right
one for pure periodicity. Today students need the definition
and examples; Lecture 2 explains the mechanism.
Source: book 3.8--3.9.
}

\teachslide{What we have proved}{2}{
\[
 \key{\text{eventually periodic}\quad\Longrightarrow\quad
 \text{quadratic irrational}}
\]
\spacer
A periodic tail satisfies a fixed-point equation.

A finite prefix changes $N$ to $MNM^{-1}$.
\spacer
Next lecture: why every quadratic irrational repeats,
and why reduced surds repeat immediately.
\spacer
\textbf{Exit question:} Why does $t=T_I(t)$ prove nothing?
}{
Expected exit answer: the identity transformation fixes
every real number, so the resulting polynomial is $0=0$.
We needed a nonscalar transformation, supplied by the
positive lower-left entry of the period matrix.
Stop here at 50 minutes. The remaining slides are backups,
not part of the planned lecture. Suggested preparation:
review complete quotients and the fact that convergents
$p_k/q_k$ tend to the original irrational.
}

\backups

\backupslide{Composition in coordinates}{
\[
 M=\mat abcd,\qquad N=\mat efgh
\]
\[
 \begin{aligned}
 T_M(T_N(t))
 &=\frac{a(et+f)+b(gt+h)}{c(et+f)+d(gt+h)}\\[8pt]
 &=\frac{(ae+bg)t+(af+bh)}{(ce+dg)t+(cf+dh)}.
 \end{aligned}
\]
\spacer
These are exactly the four entries of $MN$.
}{
Use only if students request the general composition
calculation. The equation holds wherever the displayed
finite formulas are defined, and globally if infinity
is included in the projective convention. Source: book 3.2.
}

\backupslide{The nonscalar fixed-point lemma}{
Let $M=\mat abcd$ be an invertible rational matrix.
\[
 t=T_M(t)\quad\Longrightarrow\quad
 ct^2+(d-a)t-b=0.
\]
If all coefficients vanish, then $M=aI$.

If $c=0$ and the polynomial is nonzero, its root is rational.
\spacer
Therefore an irrational fixed point of a nonscalar rational
matrix is a quadratic irrational.
}{
This separates the elementary lemma from any special
continued-fraction properties. Note that being invertible
does not by itself exclude the identity transformation.
One needs the nonscalar assumption as well as an
irrational fixed point. Source: book Lemma 3.12.
}

\backupslide{Why an infinite simple continued fraction is irrational}{
The convergents are in lowest terms, $q_k\to\infty$, and
\[
 \left|x-\frac{p_k}{q_k}\right|
 \le\frac1{q_kq_{k+1}}.
\]
If $x=p/q$ were rational, then eventually $p_k/q_k\ne p/q$,
and
\[
 \left|\frac pq-\frac{p_k}{q_k}\right|
 \ge\frac1{qq_k}.
\]
For $q_{k+1}>q$, these inequalities contradict each other.
}{
The alternating-convergent argument produces a limit
for any infinite positive digit sequence, not just
for the expansion of a number already known to be irrational.
For a rational candidate $p/q$ in lowest terms, increasing
denominators ensure that the convergents eventually
differ from it. The integer numerator of their difference
then has absolute value at least one. A weak upper bound
is enough here. Source: book 3.7.
}
